\section[Cuando los datos desaparecen]{Cuando los datos desaparecen:\\
Limpieza, reconstrucción y análisis del comportamiento USD/MXN mediante medias móviles}

En muchos libros de estadística, ciencia de datos y aprendizaje automático, los conjuntos de datos aparecen completos y listos para analizarse. Pareciera que el investigador siempre dispone de toda la información necesaria y que los datos existen sin errores, sin huecos y sin interrupciones.

Sin embargo, en problemas reales esto rara vez ocurre. En muchas ocasiones existen fechas faltantes, errores de captura, interrupciones temporales o datos atípicos.

Esto sucede incluso en variables ampliamente monitoreadas, como el tipo de cambio USD/MXN. Aunque existen múltiples plataformas financieras, no siempre se dispone de una serie perfecta.

Por ello, antes de realizar análisis estadístico, predicción o modelado, es necesario limpiar y reconstruir los datos.

Esta práctica estudia distintas estrategias de limpieza de datos utilizando la serie histórica USD/MXN.

%===========================================================
\subsection{Fuente de datos}
%===========================================================

Fuente utilizada:

\url{https://mx.investing.com/currencies/usd-mxn-historical-data}

Periodo analizado:

\[
03/03/2025 \text{ a } 17/05/2026
\]

%===========================================================
\subsection{Objetivos del estudio}
%===========================================================

\begin{itemize}
\item Detectar fechas faltantes.
\item Detectar datos atípicos.
\item Comparar técnicas de limpieza.
\item Aplicar medias móviles de tamaño 2 a 10.
\item Comparar errores mediante MSE.
\item Analizar la distribución del error.
\end{itemize}

%===========================================================
\subsection{Variables utilizadas}
%===========================================================

\begin{table}[H]
\centering
\begin{tabular}{ll}
\toprule
Variable & Descripción\\
\midrule
Fecha & Día de observación\\
Cierre & Tipo de cambio de cierre\\
Apertura & Tipo de cambio de apertura\\
Máximo & Valor máximo diario\\
Mínimo & Valor mínimo diario\\
\% var. & Variación porcentual diaria\\
\bottomrule
\end{tabular}
\end{table}

%===========================================================
\subsection{Metodología de limpieza}
%===========================================================

%===========================================================
\subsubsection{Detección de datos faltantes}
%===========================================================

Se construye un calendario diario continuo y posteriormente se comparan las fechas observadas contra las fechas esperadas.

Esto permite identificar:

\begin{itemize}
\item fechas faltantes,
\item huecos temporales,
\item discontinuidades,
\item registros incompletos.
\end{itemize}

%===========================================================
\subsubsection{Detección de datos atípicos}
%===========================================================

Además de detectar fechas faltantes, también es importante identificar datos atípicos. Un dato atípico puede representar:

\begin{itemize}
\item un error de captura,
\item un problema de almacenamiento,
\item un cambio extraordinario del mercado,
\item o un evento financiero real.
\end{itemize}

Por ello, en esta práctica los datos atípicos se marcan para análisis, pero no se eliminan automáticamente.

%===========================================================
\paragraph{Método del rango intercuartílico (IQR)}
%===========================================================

\[
IQR = Q_3 - Q_1
\]

Un valor se considera atípico si:

\[
y_i < Q_1 - 1.5(IQR)
\]

o

\[
y_i > Q_3 + 1.5(IQR)
\]

%===========================================================
\paragraph{Método del Z-score}
%===========================================================

\[
z_i = \frac{y_i-\bar y}{s}
\]

Se considera posible dato atípico cuando:

\[
|z_i| > 3
\]

%===========================================================
\subsubsection{Técnicas de reconstrucción de datos faltantes}
%===========================================================

%===========================================================
\paragraph{Eliminación (listwise deletion)}
%===========================================================

Elimina observaciones incompletas.

\textbf{Ventajas:}
\begin{itemize}
\item no inventa información,
\item simple de implementar.
\end{itemize}

\textbf{Desventajas:}
\begin{itemize}
\item reduce tamaño de muestra,
\item rompe continuidad temporal.
\end{itemize}

%===========================================================
\paragraph{Interpolación lineal}
%===========================================================

\[
y_t =
y_{t-1}
+
\frac{y_{t+1}-y_{t-1}}{2}
\]

%===========================================================
\paragraph{Promedio local (ventana simétrica de 6 vecinos)}
%===========================================================

Se utilizan tres datos previos y tres posteriores:

\[
y_t =
\frac{
y_{t-3}+y_{t-2}+y_{t-1}
+y_{t+1}+y_{t+2}+y_{t+3}
}{6}
\]

%===========================================================
\paragraph{Promedio mensual}
%===========================================================

\[
\bar y_m =
\frac{1}{n_m}
\sum_{i=1}^{n_m} y_i
\]

%===========================================================
\subsection{Análisis mediante medias móviles}
%===========================================================

Después de limpiar los datos se aplican medias móviles:

\[
n=2,3,4,\dots,10
\]

La media móvil se define como:

\[
MM_n =
\frac{1}{n}
\sum_{i=1}^{n} y_i
\]

%===========================================================
\subsection{Métricas de evaluación}
%===========================================================

%===========================================================
\subsubsection{Error cuadrático medio (MSE)}
%===========================================================

\[
MSE =
\frac{1}{n}
\sum_{i=1}^{n}
(y_i-\hat y_i)^2
\]

donde:

\begin{itemize}
\item \(y_i\): valor observado,
\item \(\hat y_i\): valor estimado.
\end{itemize}

%===========================================================
\subsubsection{Distribución del error}
%===========================================================

El error se define como:

\[
e_i = y_i - \hat y_i
\]

Para visualizar la distribución del error se utilizan diagramas de violín.

Estos permiten observar:

\begin{itemize}
\item dispersión,
\item concentración,
\item asimetría,
\item presencia de valores extremos.
\end{itemize}

%===========================================================
\subsection{Resultados esperados}
%===========================================================

\begin{itemize}
\item Fechas faltantes detectadas.
\item Datos atípicos identificados.
\item Series reconstruidas.
\item Medias móviles calculadas.
\item Comparación de MSE entre técnicas.
\item Diagramas de violín generados.
\item Mejor técnica de limpieza identificada.
\item Mejor tamaño de ventana seleccionado.
\end{itemize}

%===========================================================
\subsection{Conclusiones}
%===========================================================

La limpieza de datos modifica directamente los resultados del análisis. En series financieras, la manera en que se reconstruyen datos faltantes puede alterar tendencias, errores y conclusiones.

Por ello, antes de aplicar modelos estadísticos o predictivos, es necesario comprender qué técnica de limpieza se utiliza y cómo afecta los resultados.

%===========================================================
%===========================================================

%===========================================================
\subsection{Anexo A: Pseudocódigo del procedimiento}
%===========================================================

\begin{verbatim}
INICIO

Leer datos USD/MXN

Construir calendario completo

Detectar:
    Fechas faltantes
    Datos atípicos

Aplicar técnicas de reconstruccion:
    Eliminacion
    Interpolacion lineal
    Promedio local
    Promedio mensual

Calcular:
    Medias moviles (n = 2 a 10)
    MSE para cada media movil
    Distribucion del error

Comparar resultados entre tecnicas

FIN
\end{verbatim}

%===========================================================
\subsection{Anexo B: Muestra de datos utilizados}
%===========================================================

\begin{verbatim}
15/05/2026 17.3400 17.2215 17.4085 17.2090 +0.69%
14/05/2026 17.2215 17.1775 17.2490 17.1575 +0.26%
13/05/2026 17.1775 17.2280 17.2582 17.1629 -0.27%
12/05/2026 17.2239 17.2025 17.2850 17.1805 +0.15%
11/05/2026 17.1975 17.1915 17.2391 17.1587 -0.11%
10/05/2026 17.2171 17.1810 17.2525 17.2090 +0.20%
...
03/03/2025 20.6860 20.5265 20.7505 20.3722 +0.75%
\end{verbatim}

%===========================================================
\subsection{Anexo C: Implementación en Python}
%===========================================================

\begin{lstlisting}[language=Python]
import pandas as pd
import numpy as np
import matplotlib.pyplot as plt
import re

raw_text = """
PEGAR AQUI LOS DATOS COMPLETOS
"""

patron = re.compile(
    r"(\d{2}/\d{2}/\d{4})\s+"
    r"([\d.]+)\s+"
    r"([\d.]+)\s+"
    r"([\d.]+)\s+"
    r"([\d.]+)\s+"
    r"([+-]?[\d.]+)%"
)

filas = []

for linea in raw_text.splitlines():

    m = patron.search(linea)

    if m:
        filas.append(m.groups())

df = pd.DataFrame(
    filas,
    columns=[
        "Fecha",
        "Cierre",
        "Apertura",
        "Maximo",
        "Minimo",
        "Var_pct"
    ]
)

df["Fecha"] = pd.to_datetime(
    df["Fecha"],
    format="%d/%m/%Y"
)

for col in [
    "Cierre",
    "Apertura",
    "Maximo",
    "Minimo",
    "Var_pct"
]:
    df[col] = pd.to_numeric(
        df[col],
        errors="coerce"
    )

# DETECCION DE ATIPICOS CON IQR

def detectar_atipicos_iqr(data, columna):

    q1 = data[columna].quantile(0.25)
    q3 = data[columna].quantile(0.75)

    iqr = q3 - q1

    li = q1 - 1.5 * iqr
    ls = q3 + 1.5 * iqr

    return (
        (data[columna] < li) |
        (data[columna] > ls)
    )

df["Atipico"] = detectar_atipicos_iqr(
    df,
    "Cierre"
)

# CREACION DE CALENDARIO COMPLETO

calendario = pd.DataFrame({
    "Fecha": pd.date_range(
        df["Fecha"].min(),
        df["Fecha"].max(),
        freq="D"
    )
})

df_completo = calendario.merge(
    df,
    on="Fecha",
    how="left"
)

# INTERPOLACION LINEAL

columnas_numericas = [
    "Cierre",
    "Apertura",
    "Maximo",
    "Minimo",
    "Var_pct"
]

df_interpolacion = df_completo.copy()

df_interpolacion[columnas_numericas] = (
    df_interpolacion[columnas_numericas]
    .interpolate(method="linear")
)

# CALCULO DE MEDIAS MOVILES

for n in range(2,11):

    df_interpolacion[f"MM_{n}"] = (
        df_interpolacion["Cierre"]
        .rolling(window=n, min_periods=1)
        .mean()
    )

# FUNCION PARA CALCULAR MSE

def calcular_mse(real, predicho):

    return np.mean(
        (real - predicho)**2
    )

# TABLA DE ERRORES POR VENTANA

tabla_errores = []

for n in range(2,11):

    mse = calcular_mse(
        df_interpolacion["Cierre"],
        df_interpolacion[f"MM_{n}"]
    )

    tabla_errores.append({
        "n": n,
        "MSE": mse
    })

tabla_errores = pd.DataFrame(
    tabla_errores
)

print(tabla_errores)

# DIAGRAMA DE VIOLIN DE ERRORES

errores = []

for n in range(2,11):

    temp = (
        df_interpolacion["Cierre"]
        -
        df_interpolacion[f"MM_{n}"]
    )

    errores.append(
        temp.dropna()
    )

plt.figure(figsize=(12,6))

plt.violinplot(
    errores,
    showmeans=True,
    showmedians=True
)

plt.xticks(
    range(1,10),
    [f"MM_{n}" for n in range(2,11)]
)

plt.title(
    "Distribucion del error por media movil"
)

plt.grid(True)

plt.show()
\end{lstlisting}
